
\documentclass{article}
\usepackage{colortbl}
\usepackage{makecell}
\usepackage{multirow}
\usepackage{supertabular}

\begin{document}

\newcounter{utterance}

\centering \large Interaction Transcript for game `eqbench', experiment `eq\_bench', episode 124 with reflexion{-}Qwen3.5{-}9B.
\vspace{24pt}

{ \footnotesize  \setcounter{utterance}{1}
\setlength{\tabcolsep}{0pt}
\begin{supertabular}{c@{$\;$}|p{.15\linewidth}@{}p{.15\linewidth}p{.15\linewidth}p{.15\linewidth}p{.15\linewidth}p{.15\linewidth}}
    \# & \multicolumn{2}{c}{Player} && \multicolumn{2}{c}{Game Master} \\
    \hline

    \theutterance \stepcounter{utterance}  
    & & & \multicolumn{4}{p{0.6\linewidth}}{
        \cellcolor[rgb]{0.9,0.9,0.9}{
            \makecell[{{p{\linewidth}}}]{
                \texttt{\tiny{[P1$\langle$GM]}}
                \texttt{You have attempted this same task before and did not succeed. These are the plans you wrote after those attempts.} \\
\texttt{Trial \#1: New plan: I failed because I assigned a high score to "Contemptuous" (8) and "Insecure" (6) while underestimating "Offended" (2). In this specific dialogue, Argyle's response ("Like the respect you show...") is a direct, personal attack on Serepta's character and actions, which is the primary trigger for feeling offended. While Serepta might feel insecure about her control or contemptuous of Argyle's attitude, the immediate and most intense reaction to being accused of lacking respect and being controlling is offense. Next time, I will prioritize "Offended" as the highest score (likely 9 or 10) since the dialogue centers on a conflict of respect and boundaries, and I will lower the scores for secondary emotions to reflect that they are reactions to the primary feeling of being disrespected.} \\
\\ 
\texttt{Trial \#2: New plan: I failed because I assigned a high score to "Contemptuous" (7) and "Insecure" (4), underestimating the intensity of "Offended" (9) relative to the other options. While Serepta is certainly feeling insecure about her lack of control and contemptuous of Argyle's attitude, the dialogue explicitly frames the conflict around "respect" and "boundaries." Argyle's response ("Like the respect you show...") is a direct, personal accusation that triggers a defensive and hurt reaction, making "Offended" the dominant and most intense emotion. Next time, I will assign the maximum score (10) to "Offended" as the primary reaction to the direct attack on her character, and significantly lower the scores for "Contemptuous" (2) and "Insecure" (3) to reflect that they are secondary to the immediate feeling of being disrespected. Disinterested will remain 0.} \\
\\ 
\texttt{Trial \#3: New plan: I failed because my previous attempts underestimated the intensity of "Offended" relative to "Contemptuous" and "Insecure," likely due to misinterpreting Argyle's defensive retort as a sign of Serepta's own insecurity or arrogance rather than a direct trigger for her feeling disrespected. Argyle's accusation ("Like the respect you show...") is a personal attack on Serepta's character, which is the most potent catalyst for feeling offended in this context. While Serepta may feel insecure about her lack of control or contemptuous of Argyle's messiness, these are secondary to the immediate sting of being accused of lacking respect. Next time, I will assign the maximum score (10) to "Offended" as the primary and most intense reaction to the direct personal accusation, while significantly lowering "Contemptuous" to 2 and "Insecure" to 3 to reflect their status as secondary, underlying feelings rather than the dominant response. "Disinterested" will remain 0.} \\
\\ 
\\ 
\\ 
\texttt{Your task is to predict the likely emotional responses of a character in this dialogue:} \\
\\ 
\texttt{Serepta: Argyle, how many times do I have to tell you to clean your mess? } \\
\texttt{Argyle: Serepta, it's just a few sketches scattered around. It's not like I'm hoarding garbage.} \\
\texttt{Serepta: It's not about the sketches, Argyle. It's about respect. This is our shared space.} \\
\texttt{Argyle: Respect, huh? Like the respect you show when you rearrange my things without asking?} \\
\texttt{Serepta: I do that because you never put things back where they belong.} \\
\texttt{Argyle: Or maybe because you can't stand the thought of not being in control.} \\
\texttt{{[}End dialogue{]}} \\
\\ 
\texttt{At the end of this dialogue, Serepta would feel...} \\
\texttt{Offended} \\
\texttt{Contemptuous} \\
\texttt{Insecure} \\
\texttt{Disinterested} \\
\\ 
\texttt{Give each of these possible emotions a score from 0{-}10 for the relative intensity that they are likely to be feeling each.} \\
\\ 
\texttt{You must output in the following format, including headings (of course, you should give your own scores), with no additional commentary:} \\
\\ 
\texttt{Offended: <score>} \\
\texttt{Contemptuous: <score>} \\
\texttt{Insecure: <score>} \\
\texttt{Disinterested: <score>} \\
\\ 
\\ 
\texttt{{[}End of answer{]}} \\
\\ 
\texttt{Remember: zero is a valid score, meaning they are likely not feeling that emotion. You must score at least one emotion > 0.} \\
\\ 
\texttt{Your answer:} \\
            }
        }
    }
     \\ \\

    \theutterance \stepcounter{utterance}  
    & \multicolumn{4}{p{0.6\linewidth}}{
        \cellcolor[rgb]{0.561,0.737,0.561}{
            \makecell[{{p{\linewidth}}}]{
                \texttt{\tiny{[P1$\rangle$GM]}}
                \texttt{Offended: 10} \\
\texttt{Contemptuous: 2} \\
\texttt{Insecure: 3} \\
\texttt{Disinterested: 0} \\
            }
        }
    }
    & & \\ \\

    \theutterance \stepcounter{utterance}  
    & & & \multicolumn{2}{p{0.3\linewidth}}{
        \cellcolor[rgb]{0.9,0.9,0.9}{
            \makecell[{{p{\linewidth}}}]{
                \texttt{\tiny{[GM$|$GM]}}
                \texttt{Offended: 10} \\
\texttt{Contemptuous: 2} \\
\texttt{Insecure: 3} \\
\texttt{Disinterested: 0} \\
            }
        }
    }
    & & \\ \\

    \theutterance \stepcounter{utterance}  
    & & & \multicolumn{2}{p{0.3\linewidth}}{
        \cellcolor[rgb]{0.9,0.9,0.9}{
            \makecell[{{p{\linewidth}}}]{
                \texttt{\tiny{[GM$|$GM]}}
                \texttt{Offended: 8} \\
\texttt{Contemptuous: 6} \\
\texttt{Insecure: 0} \\
\texttt{Disinterested: 0} \\
            }
        }
    }
    & & \\ \\

    \theutterance \stepcounter{utterance}  
    & & & \multicolumn{2}{p{0.3\linewidth}}{
        \cellcolor[rgb]{0.9,0.9,0.9}{
            \makecell[{{p{\linewidth}}}]{
                \texttt{\tiny{[GM$|$GM]}}
                \texttt{game\_result = LOSE} \\
            }
        }
    }
    & & \\ \\

    \theutterance \stepcounter{utterance}  
    & & \multicolumn{4}{p{0.6\linewidth}}{
        \cellcolor[rgb]{0.95,0.95,0.95}{
            \makecell[{{p{\linewidth}}}]{
                \texttt{\tiny{[P1$\leftrightarrow$P]}}
                \texttt{Post{-}episode Reflexion (not sent to GM):} \\
\texttt{New plan: I failed because I assigned a high score to "Contemptuous" and "Insecure" while underestimating the intensity of "Offended." In this specific dialogue, Argyle's response ("Like the respect you show...") is a direct, personal attack on Serepta's character and actions, which is the primary trigger for feeling offended. While Serepta might feel insecure about her control or contemptuous of Argyle's attitude, the immediate and most intense reaction to being accused of lacking respect and being controlling is offense. Next time, I will prioritize "Offended" as the highest score (10) since the dialogue centers on a conflict of respect and boundaries, and I will lower the scores for secondary emotions to reflect that they are reactions to the primary feeling of being disrespected.} \\
            }
        }
    }
    & \\ \\

\end{supertabular}
}

\end{document}
